\chapter{Identity: What Survives Mutation}\label{ch:identity}

\NewV\ \emph{This chapter does not exist in the V0 book. It is the centre of
V2 and everything after it depends on it.}

\begin{workplan}{\NewV}
\wpgoal{Make ``\textsc{tatiana} stays itself while it changes'' into a number
that can be printed and a theorem that can fail. Nothing in this chapter may
rest on the word ``identity'' being intuitively clear.}

\wpsource{\texttt{DOCS/TATIANA\_V1\_IDENTITY\_MAP.md}~\S3 (the three levels),
\S4 (which phase conserves what);
\texttt{MATH\_TOOLBOX.md}~II.2 (the $\mathbb{K}/W$ split and the consolidation
formula), III.3 (Construction~6 fixes the growth increment);
\texttt{TATIANA-V2/MATH/phase0\_invariant\_dimension.md} (the first
measurement, already run --- see step~5).}

\wpsteps{
  \item \textbf{State the split first.} Identity is carried by the
    crystallized store $\mathbb{K}$, not by the working complex $W$; $W$ may
    be destroyed at the end of every session. Give the reader the reason this
    is the right place for identity to live \emph{before} any formula: a
    memory that is destroyed nightly cannot be what makes you you.
  \item \textbf{Level 0 --- the exact invariant.} Define
    $\mathrm{K}(\mathbb{K}) = (\dim \Fs_{\mathbb{K}}(\sigma))_\sigma$, the
    dimension vector, a $K_0$ class. Prove that the consolidation rule
    $R^{\mathbb{K}} \leftarrow \Pi_{O(d)}(R^{\mathbb{K}} + \gamma(\nu)(R^W -
    R^{\mathbb{K}}))$ conserves it \emph{exactly}. The proof is short and must
    be given in full: a retraction onto $O(d)$ cannot change a dimension,
    molding cannot, state relaxation cannot.
  \item \textbf{The one thing that moves it.} Growth, and only growth, and by
    a \emph{derived} amount: $\dim \Fs(w) = \dim \Hzero(\gamma; \Fs|_\gamma)$
    from Construction~6 (Chapter~\ref{ch:cone}). Nothing is hand-set. State
    the theorem shape explicitly, because it is the shape V2 is built to
    have: \emph{two of three timescales conserve the invariant exactly; the
    third changes it computably.}
  \item \textbf{Level 1 --- the projective invariant.} Level 0 is coarse: it
    cannot distinguish two stores with the same dimensions and different maps.
    State the finer claim, that relative structure is preserved while absolute
    magnitude is not, $\Fs_{\mathbb{K}} \sim \lambda\Fs_{\mathbb{K}}$. Then
    make the observation that earns it: $\Pi_{O(d)}$ is the polar retraction
    $R \mapsto UV^{\mathsf T}$, which discards scale and keeps direction ---
    written into the design months earlier as a gauge fixing, and only later
    recognised as the identity-preservation mechanism. Say plainly that the
    coincidence is the reason to trust it, and that this is
    \Contested\ until step~5's measurement is extended to cover it directly.
  \item \textbf{Report the measurement that exists.} Phase~0 has been run.
    Reproduce its result here in full, including its controls: on a $3$-cycle
    and on a $b_1=0$ path, pre-existing $\mathrm{K}(\mathbb{K})$ is untouched
    across 500 wake/sleep ticks; the single growth event's size is $k_w = 1$
    confirmed by two independent routes (holonomy $\dim\ker(H-\Id)$ and
    cochain-space $\dim\ker(\cob^0(\cob^0)^{\mathsf T})$); $Q(t)$ moves while
    dimensions stay pinned. Tag \Measured.
  \item \textbf{State what the measurement does not cover, in the same
    breath.} Sleep-phase pruning is unimplemented, so the ``may decrease''
    half of the invariant table is untested and the kill-test is
    \emph{survived, not passed}. No $2$-cells; $\gamma(\nu)$ constant;
    post-growth dynamics unwired; dense rather than Householder maps;
    synthetic input only. This list is not an appendix item --- it goes
    beside the result.
  \item \textbf{Level 2 --- declare it out of scope and say why.} The full
    isomorphism class of $\Fs_{\mathbb{K}}$: a cellular sheaf on a graph is a
    representation of the incidence quiver, and anything richer than a single
    cycle is of wild representation type. Krull--Schmidt makes the
    decomposition into indecomposables unique but not computable. V2 reports
    levels 0 and 1 and does not chase level 2.
}

\wpdone{A reader can state what $\mathrm{K}(\mathbb{K})$ is, prove
unaided that consolidation conserves it, say exactly how much growth changes
it and why that number is not a choice, and list what the existing
measurement does not cover.}
\end{workplan}

\begin{workplan}{\NewV}
\wpgoal{\S: The open question this chapter raises and cannot yet close ---
whether molding can erode the obstruction that growth depends on.}

\wpsource{Raised by the Phase~0 experiment; see
\texttt{TATIANA-V2/MATH/phase0\_invariant\_dimension.md}~\S4 and the logbook
entry for 2026-08-14.}

\wpsteps{
  \item State the concern precisely. Molding changes restriction maps; the
    growth address is a harmonic class determined \emph{by} those maps. In the
    experiment as run, molding is bounded to one nudge per tick before the
    same tick's stall check, so the question never arises. Unbounded, it might.
  \item Say why it matters: if slow molding can dissolve a harmonic
    obstruction, the three-layer handoff of Chapter~\ref{ch:twophase} is not
    a handoff but a race, and ``growth fires when molding stalls'' would be
    a statement about the schedule rather than about the geometry.
  \item Do not answer it here. Record it as an open problem with a stated
    experiment (vary the molding budget per tick; measure whether the harmonic
    component's norm decays to below the stall tolerance without growth ever
    firing) and cross-reference Chapter~\ref{ch:open}.
}

\wpdone{The question is stated sharply enough that someone else could run
the experiment that settles it.}
\end{workplan}
